\documentclass[10pt,a4paper]{amsart}
\usepackage[margin=19mm]{geometry}
\usepackage{amsmath,amssymb,mathtools}
\usepackage[scr=rsfs,cal=euler]{mathalfa}
\usepackage{xfrac}
\usepackage[unicode,hidelinks,bookmarks=false]{hyperref}
\newcommand{\ie}{i.e.\ }
\newcommand{\cf}{cf.\ }
\newcommand{\resp}{respectively\ }
\newcommand{\Subsection}{\subsection}
\providecommand{\nobreakdash}{\nobreak\hbox{-}}
\setlength{\emergencystretch}{2em}
\multlinegap0pt
\usepackage{bm}
\usepackage{bbm}
\usepackage{dsfont}
\usepackage{xspace}
\usepackage{cancel}
\usepackage{mathtools}
\usepackage{amsthm}
\makeatletter
\@ifundefined{prop}{\newtheorem{prop}{Proposition}}{}
\makeatother
\makeatletter
\DeclareMathOperator{\Ad}{Ad}
\DeclareMathOperator{\Supp}{Supp}
\DeclareMathOperator{\cO}{\mathcal{O}}
\expandafter\ifx\csname num\endcsname\relax
\newenvironment{num}
{\setcounter{keepeqno}{\value{equation}}%
	\begin{list}{(\theequation)}{\usecounter{equation}}%
		\setcounter{equation}{\value{keepeqno}}}
	{\end{list}}
\fi
\def\oN{\overline{N}}
\def\ou{\overline{u}}
\expandafter\ifx\csname paragr\endcsname\relax
\newenvironment{paragr}[1][]{\refstepcounter{subsubsection} \noindent \textbf{\thesubsubsection . \ #1}}{\medskip}
\fi
\DeclareMathOperator{\rI}{\mathrm{I}}
\DeclareMathOperator{\rO}{\mathrm{O}}
\def\tM{\widetilde{M}}
\def\tm{\widetilde{m}}
\makeatother
\title[The Hiraga-Ichino-Ikeda conjecture on formal degrees for cla]{The Hiraga-Ichino-Ikeda conjecture on formal degrees for classical groups}
\author{Rapha\"el Beuzart-Plessis}
\date{Source: arXiv:2508.08470v1 (CC BY 4.0)}
\begin{document}
\maketitle

\noindent\emph{Source and licence.} The Hiraga-Ichino-Ikeda conjecture on formal degrees for classical groups. Version arXiv:2508.08470v1 (CC BY 4.0), distributed under the Creative Commons Attribution 4.0 International licence, as recorded in the arXiv licence metadata for this exact version and shown on the abstract page https://arxiv.org/abs/2508.08470v1. Full rights notice: licenses/arxiv-2508.08470-CC-BY-4.0.txt.\par\smallskip

\noindent\textit{Editorial scope.} Counted unit: the Prop whose source label is \texttt{prop} in the article (source0.tex lines 815--823, source label \texttt{prop}), delivered together with the article's own complete original proof of it (source0.tex lines 825--843, one \texttt{proof} environment of 19 lines). No mathematics is written, repaired or re-derived: statement and proof are the article's own text line for line. Required context retained uncounted: eq tildedelta (lines 780--782); eq2 (lines 786--788); eq tf (lines 810--812). Every definition, display and label of those lines is retained unchanged. Omitted from this package and not redistributed: 1--779, 783--785, 789--809, 813--814, 824--824, 844--1491. Nothing inside a retained excerpt is omitted or altered: every retained body is delivered exactly as the article's lines are written, so no omission receipt of any kind accompanies this package. Citations: the retained excerpts carry no citation command at all, so no bibliography entry of the article is transcribed and no reference is redistributed. Reference adaptation: none was needed and none was made; every reference inside the delivered unit resolves inside it, so no unresolved reference or citation is produced. Printed slips kept literal: no printed word, symbol, hypothesis or number of the counted statement or of its proof was altered. Layout and class: the article's own class is \texttt{article} with packages including amsmath, mathrsfs, amsfonts, amsthm, mathtools, stmaryrd, amssymb, fullpage, hyperref, epic, eepic, babel, xy, enumerate; the wrapper is the lane's pinned \texttt{amsart} editorial wrapper at 10pt on A4, which restates the theorem-like environments the retained text needs and adds the article's own macro definitions that the retained text needs (\texttt{\textbackslash Ad}, \texttt{\textbackslash Supp}, \texttt{\textbackslash cO}, \texttt{\textbackslash oN}, \texttt{\textbackslash ou}, \texttt{\textbackslash rI}, \texttt{\textbackslash rO}, \texttt{\textbackslash tM}, \texttt{\textbackslash tm}), with extra wrapper packages (bm, bbm, dsfont, xspace, cancel, mathtools). The delivered numbering is the unit's own. Assets: no retained span contains a figure environment, a graphics-inclusion command, a PDF-page inclusion, a table or a TikZ picture; no image file is copied into this package, the package declares no asset dependency and no image file is redistributed with it. Licence: version arXiv:2508.08470v1 of this article is distributed under the Creative Commons Attribution 4.0 International licence, as recorded in the arXiv licence metadata for this exact version and shown on the abstract page https://arxiv.org/abs/2508.08470v1. A full rights notice is delivered at licenses/arxiv-2508.08470-CC-BY-4.0.txt. Characters: every retained excerpt is pure ASCII, so the delivered engine, XeLaTeX with its default Latin Modern text font, reports no missing character.
	\begin{equation}\label{eq tildedelta}
		\displaystyle \widetilde{\delta}(m_1\tm m_2)=\delta(m_1)^{-1}\widetilde{\delta}(\tm)\delta(m_2),\;\; \mbox{ for every } (m_1,\tm,m_2)\in M\times \tM\times M.
	\end{equation}

	\begin{equation}\label{eq2}
		\displaystyle \int_{\oN'} f(\tm(u)) du= C_0 \cdot \rO(\gamma_{-1},\widetilde{\delta}^{1/2}f)
	\end{equation}

	\begin{equation}\label{eq tf}
	\displaystyle f=\widetilde{\delta}^{1/2} \cdot \Phi^\vee\mid_{\tM},
\end{equation}

	\begin{prop}\label{prop}
		The integral
		\begin{equation}\label{eq4}
			\displaystyle \int_{\oN} \Phi_s(\ou) d\ou
		\end{equation}
		converges for $\Re(s)>0$ and moreover we have the identity
		$$\displaystyle \lim\limits_{s\to 0^+} \gamma(s,\mathbf{1}_F,\psi) \int_{\oN} \Phi_s(\ou)d\ou =C_1 \cdot \rI_{\chi}(f)$$
		where $C_1>0$ is an absolute constant (depending on our choice of Haar measures). 
	\end{prop}

	\begin{proof}
		Let $A^2$ denote the subgroup of squares in $A$. We have the integration formula\footnote{Indeed, it suffices to check that it holds for the characteristic function $\phi=\mathbf{1}_{\cO_F^\times}$; note that $2\lvert 2\rvert^{-1}$ is precisely the cardinality of the group of square classes $\cO_F^\times/\cO_F^{\times,2}$.}
		$$\displaystyle \int_A \phi(a)da=\frac{\lvert 2\rvert}{2}\sum_{t\in A/A^2} \int_A \phi(ta^2) da,\;\; \phi\in C_c^\infty(A).$$
		Thus, by definition of the convenient section $s\mapsto \Phi_s$, and using the fact that $\chi$ is quadratic, we have (ignoring for the moment convergence issues)
		\begin{align}\label{eq3}
			\displaystyle \int_{\oN} \Phi_s(\ou) d\ou & =\int_{\oN\times A} \Phi(a\ou) \chi(a)\delta(a)^{-1/2} \lvert a\rvert^{-ds} dad\ou \\
			\nonumber & = \frac{\lvert 2\rvert}{2} \sum_{t\in A/A^2} \delta(t)^{-1/2}\chi(t) \lvert t\rvert^{-ds}\int_{\oN\times A} \Phi(ta^2\ou)\delta(a)^{-1} \lvert a\rvert^{-2ds} dad\ou \\
			\nonumber & = \frac{\lvert 2\rvert}{2}  \sum_{t\in A/A^2} \delta(t)^{-1/2}\chi(t)  \lvert t\rvert^{-ds}\int_{\oN\times A} \Phi^t(a\ou a) \lvert a\rvert^{-2ds} dad\ou
		\end{align}
		where we have set $\Phi^t(g)=\Phi(tg)$ for any $g\in G$ and the last equality follows from the change of variable $\ou\mapsto a^{-1}\ou a$. Fix $t\in A$ and let $u\in \oN'$. There exists a unique $u(\ou)\in N$ such that $\ou\in N \tm(u)u(\ou)$ and we have
		$$\displaystyle \int_A \Phi^t(a\ou a) \lvert a\rvert^{-2ds} da=\int_A \Phi^t(\tm(\ou) a u(\ou)a^{-1}) \lvert a\rvert^{2ds}da.$$
		Note that the function $a\in F^\times \mapsto \Phi^t(\tm(\ou) au(\ou)a^{-1})$ vanishes for $a\in F^\times$ sufficiently large (by compactness of the support of $f$ and because $u(\ou)\neq 1$) and is constant equal to $\Phi^t(\tm(u))$ in a neighborhood of $0$ (by smoothness of $\Phi$) i.e. it extends to a function in $C_c^\infty(F)$. Thus, by Tate's thesis the above integral converges for $\Re(s)>0$ and is equivalent to $\gamma(2ds,\mathbf{1}_F,\psi)^{-1}\Phi^t(\tm(u))$ when $s\to 0$. Moreover, we claim that all of this happens uniformly in $\ou$. More precisely, as $\Supp(\Phi^t)\subset G'=N \tM N$, we may find a compact subset $\Omega\subset \tM$ such that $\Supp(\Phi^t)\subset N \Omega N$. It follows that the function $a\mapsto \Phi^t(a\ou a)$ vanishes identically unless $\ou$ belongs to the compact subset $\Omega'=\{u\in \oN'\mid \tm(u)\in \Omega \}$. Moreover, as the image of the map $\ou\in \Omega'\mapsto u(\ou)\in N$ is also compact (and avoiding $1\in N$), the previous reasoning shows that the functions $a\mapsto \Phi^t(a\ou a)$ ranges over a finite set as $\ou$ varies in $\Omega'$. Thus, we may write the integral over $\oN$ in the last line of \eqref{eq3} as a finite sum to deduce that \eqref{eq4} converges for $\Re(s)>0$ and
		\[\begin{aligned}
			\displaystyle \lim\limits_{s\to 0} \gamma(s,\mathbf{1}_F,\psi)\int_{\oN} \Phi_s(\ou)d\ou & = \frac{\lvert 2\rvert}{4d}\sum_{t\in A/A^2} \chi(t) \delta(t)^{-1/2} \int_{\oN} \Phi^t(\tm(\ou)) d\ou \\
			& = C_0\cdot \frac{\lvert 2\rvert}{4d} \sum_{t\in F^\times/F^{\times,2}} \chi(t) \int_{\Ad(M)\gamma_{-1}} f(t \tm) d_{-1}\tm \\
			& = C_1\cdot \sum_{t\in F^\times/F^{\times,2}} \chi(-t) O(\gamma_t,f)=C_1 \cdot \rI_{\chi}(f),
		\end{aligned}\]
		where we have set $C_1= C_0\frac{\lvert 2\rvert}{4d} $. Here, the second equality above follows from \eqref{eq2} as well as the relations \eqref{eq tildedelta} and \eqref{eq tf}, whereas the last equality is a consequence of the fact that the isomorphism $\Ad(M)\gamma_{-1}\simeq \Ad(M)\gamma_{-t}$, $\tm\mapsto t\tm$, is measure preserving.
	\end{proof}

\end{document}
